\chapter{Arbitrage systématique sur le crédit et les ETF obligataires}\label{ch:s2:systematic-credit-and-bond-etf-arbitrage}

Un ETF obligataire se négocie à chaque seconde, alors que la plupart de ses obligations ne se négocient pas pendant des jours. Lorsque les deux divergent, les participants autorisés créent ou rachètent des actions,
et la décote indique à quel point les obligations sont réellement illiquides. En mars 2020, Haddad, Moreira et Muir ont établi que les ETF obligataires liquides se négociaient avec de fortes décotes par rapport à
leur valeur liquidative, car les investisseurs vendaient ce qu'ils pouvaient pour lever des liquidités. Le même marché possède des facteurs systématiques. Houweling et van Zundert ont mis en évidence des primes de taille,
de faible risque, de valeur et de momentum dans les obligations d'entreprise, avec de faibles corrélations entre elles. Sur le marché obligataire synthétique de ce chapitre, le faible risque rapporte un
ratio de Sharpe de 0.95, la valeur 0.46 et le momentum 0.14. Dans un stress introduit à dessein, le cours de l'ETF se situe 5.0~\% sous sa valeur liquidative obsolète et 3.0~\% sous la valeur réelle des
obligations, et un participant autorisé peut en tirer profit pendant 28 jours. La construction est \texttt{firm.syscredit}.

\section{Facteurs de crédit}

\begin{definition}[Facteur crédit]\label{def:s2:systematic-credit-and-bond-etf-arbitrage:factor}
Un \emph{facteur crédit}\index{facteur crédit} est une caractéristique des obligations d'entreprise (spread par rapport à des pairs de notation et de duration similaires, variation récente du spread,
duration ou notation, taille de l'émetteur) qui les répartit en portefeuilles dont les rendements excédentaires par rapport aux Treasuries diffèrent durablement, après ajustement
pour leur spread duration.
\end{definition}

Houweling et van Zundert ont établi que les portefeuilles de taille, de faible risque, de valeur et de momentum ont produit des alphas économiquement et statistiquement
significatifs sur le marché des obligations d'entreprise. Ils ont tenu après les coûts de transaction et au sein des obligations liquides, et se sont bien combinés parce que leurs corrélations étaient faibles.
Mesurer des facteurs sur les obligations présente deux difficultés que n'ont pas les actions. Le rendement d'une obligation tient surtout à sa spread duration multipliée par la
variation de son spread, si bien que chaque signal doit être comparé à des obligations de duration voisine. Et les cours des obligations sont obsolètes, de sorte qu'un signal lu sur les cotations de la semaine passée peut déjà être
dépassé.

\texttt{firm.syscredit} construit 300 obligations, dont les durations vont de un à douze ans, sur une période de huit ans
(le \cref{lst:s2:systematic-credit-and-bond-etf-arbitrage:factors}). Chaque spread évolue sous l'effet d'un facteur de marché, d'une tendance propre à l'émetteur qui persiste pendant des mois,
un écart à la fair value qui se résorbe en un an, et du bruit. Les rendements excédentaires attendus croissent avec la racine carrée de la duration, si bien que les obligations courtes et sûres
rapportent davantage par unité de risque. Chaque portefeuille factoriel achète le cinquième supérieur et vend le cinquième inférieur des obligations selon son signal, en pondérant de façon que chaque côté ait une
spread duration comparable. La valeur est le spread au-dessus de la droite de spread duration ajustée, diminué de sa variation sur trois mois, ce qui évite que la valeur ne se réduise à une liste de
perdants récents. Le momentum est le resserrement du spread sur six mois. Le faible risque achète des obligations courtes et vend les longues, à risque égal.

\begin{center}
\small
\begin{tabular}{@{}lcc@{}}
\toprule
huit ans, rééquilibrage mensuel & ratio de Sharpe & corrélation avec la valeur\\
\midrule
valeur & 0.46 & 1\\
momentum & 0.14 & $-0.09$\\
faible risque & 0.95 & $-0.03$\\
combinaison à risque égal & 0.94 &\\
\bottomrule
\end{tabular}
\end{center}

Le faible risque domine la combinaison (la \cref{fig:s2:systematic-credit-and-bond-etf-arbitrage:factors}). Le momentum est faible parce que l'écart de valeur, qui
se résorbe, pèse pour une part dans le mouvement de six mois de chaque obligation. Les deux signaux se disputent les mêmes variations de spread. Houweling et van Zundert ont constaté que les deux
étaient positifs sur données réelles. Un modèle où l'un revient vers sa moyenne et l'autre suit une tendance sur des horizons qui se recouvrent ne laisse guère de place à l'un ni à l'autre.

\begin{omfigure}
\begin{tikzpicture}
\begin{axis}[width=11cm,height=5.4cm,xlabel={\small année},ylabel={\small cumulé (\%, à 10~\% de volatilité)},tick label style={font=\footnotesize},
  xmin=0,xmax=7.5,ymin=-20,ymax=85,grid=major,grid style={gray!25},table/col sep=comma,
  legend style={font=\scriptsize,at={(0.5,-0.3)},anchor=north},legend columns=4]
\addplot[omDef,thick] table[x=year,y=value]{figdata/strategies-2/19-systematic-credit-and-bond-etf-arbitrage/factors.csv};
\addplot[gray,thick] table[x=year,y=momentum]{figdata/strategies-2/19-systematic-credit-and-bond-etf-arbitrage/factors.csv};
\addplot[omThm,thick] table[x=year,y=lowrisk]{figdata/strategies-2/19-systematic-credit-and-bond-etf-arbitrage/factors.csv};
\addplot[black,very thick] table[x=year,y=combined]{figdata/strategies-2/19-systematic-credit-and-bond-etf-arbitrage/factors.csv};
\legend{valeur,momentum,faible risque,combiné}
\end{axis}
\end{tikzpicture}

\omcaption{Trois portefeuilles factoriels de crédit synthétiques et leur combinaison à risque égal, chacun ramené à une volatilité de 10~\%, sommes cumulées des rendements quotidiens.
Données~: \texttt{s2\_syscredit.market}.}\label{fig:s2:systematic-credit-and-bond-etf-arbitrage:factors}
\end{omfigure}

\section{ETF obligataires et création--rachat}

\begin{definition}[Arbitrage création--rachat d'ETF]\label{def:s2:systematic-credit-and-bond-etf-arbitrage:createredeem}
L'\emph{arbitrage création--rachat d'ETF}\index{arbitrage création--rachat d'ETF} est la transaction par laquelle un participant autorisé (livre~1, chapitre~14) remet un
panier de création pour recevoir de nouvelles actions de l'ETF lorsque celles-ci se négocient au-dessus de la valeur du panier, ou remet des actions pour recevoir le panier lorsqu'elles se négocient
en dessous, en conservant l'écart diminué de ses coûts.
\end{definition}

Pour un ETF obligataire, le panier est la partie délicate. Koont, Ma, Pastor et Zeng ont constaté que les ETF d'obligations d'entreprise choisissent des paniers de création et de rachat qui détiennent
des liquidités et seulement une partie des obligations de l'indice, d'autant plus que ces obligations sont illiquides. Ils ajustent les paniers pour corriger les déséquilibres tout en
préservant l'arbitrage.
L'inclusion dans le panier a rendu les obligations plus liquides en général, mais moins liquides lorsque les créations et les rachats étaient le plus déséquilibrés, comme pendant la crise du COVID-19.

\begin{definition}[Décote d'ETF]\label{def:s2:systematic-credit-and-bond-etf-arbitrage:discount}
Une \emph{décote d'ETF}\index{décote d'ETF} est l'écart dont le cours de marché d'un ETF est inférieur à sa valeur liquidative publiée~; pour un ETF obligataire, une partie
de celle-ci tient à la valeur liquidative obsolète, calculée à partir de cours d'obligations qui n'ont pas été négociés, et une autre partie à la pression vendeuse sur les actions que le mécanisme de création--rachat n'a pas
absorbée.
\end{definition}

\section{Transactions de portefeuille}

Une transaction de portefeuille (livre~2, chapitre~22) achète ou vend en une seule fois une liste complète d'obligations, souvent le panier d'un ETF, à un prix unique négocié. Les dealers
la valorisent par rapport à l'ETF et aux cotations des obligations, et peuvent la couvrir avec l'ETF lui-même. Lorsque l'ETF se négocie avec une décote, un client qui vend un portefeuille de
obligations paie le coût pour le dealer de les racheter ou de les détenir. Fournir de la liquidité aux transactions de portefeuille revient à être la contrepartie du participant
autorisé, sans passer par l'unité de création.

\section{Décotes en période de stress}

L'ETF synthétique détient les 300 obligations. Sa valeur liquidative utilise les derniers cours négociés des obligations~; chaque obligation ne se négocie que sur un quart des jours. Son cours de marché suit la
valeur réelle des obligations. Six ans après le début, un stress introduit à dessein élargit les spreads de marché de 150 points de base sur trois semaines, et le panier perd 9.8~\%. Pendant les mêmes
semaines, le cours de l'ETF tombe encore de 3~\% sous sa valeur réelle, tandis que les investisseurs vendent des actions de l'ETF pour obtenir des liquidités, puis se redresse
en deux mois (la \cref{fig:s2:systematic-credit-and-bond-etf-arbitrage:etf}).

En temps normal, la décote s'établit en moyenne à 2.7 points de base, avec un écart-type de 21.6. Au pire du stress, quatorze jours de négociation après son début, l'ETF se négociait
504 points de base sous la valeur liquidative. Sur cet écart, la valeur liquidative obsolète se situait à 215 points de base au-dessus de la valeur réelle des obligations, et le cours à 300 points de base en dessous.
Seule la seconde partie relève de l'arbitrage. Un participant autorisé qui achète des actions, les rachète et vend les obligations à leur valeur réelle paie un coût de
1.5~\% pendant le stress (0.3~\% en temps normal). La transaction a été rentable pendant 28 jours, pour un gain pouvant atteindre 1.63~\% de la valeur rachetée, et aucun jour en dehors du stress.

\begin{omfigure}
\begin{tikzpicture}
\begin{axis}[width=11cm,height=5.4cm,xlabel={\small jours de négociation depuis le début du stress},ylabel={\small valeur (100 la veille)},
  tick label style={font=\footnotesize},xmin=-10,xmax=64,ymin=86,ymax=102,grid=major,grid style={gray!25},table/col sep=comma,
  legend style={font=\scriptsize,at={(0.97,0.97)},anchor=north east}]
\addplot[black,thick] table[x=day,y=true]{figdata/strategies-2/19-systematic-credit-and-bond-etf-arbitrage/etf.csv};
\addplot[omDef,thick,dashed] table[x=day,y=nav]{figdata/strategies-2/19-systematic-credit-and-bond-etf-arbitrage/etf.csv};
\addplot[omThm,thick] table[x=day,y=price]{figdata/strategies-2/19-systematic-credit-and-bond-etf-arbitrage/etf.csv};
\legend{valeur réelle des obligations,valeur liquidative sur cours obsolètes,cours de l'ETF}
\end{axis}
\end{tikzpicture}

\omcaption{L'ETF obligataire synthétique lors d'un stress de trois semaines~: la valeur réelle de ses 300 obligations, la valeur liquidative calculée à partir de leurs derniers cours échangés, et le
cours de marché de l'ETF, qui passe sous les deux. Données~: \texttt{s2\_syscredit.market}.}\label{fig:s2:systematic-credit-and-bond-etf-arbitrage:etf}
\end{omfigure}

Haddad, Moreira et Muir interprètent de même les décotes de mars 2020. Les investisseurs ont cherché à vendre leurs positions les plus sûres et les plus liquides afin de se procurer des liquidités, et les
décotes ont été plus fortes pour les ETF sur bons du Trésor, sur obligations municipales et sur titres investment grade que pour ceux sur titres high yield. Les perturbations se sont résorbées presque aussi vite que
elles étaient apparues, après que la Réserve fédérale a annoncé des achats d'obligations d'entreprises les 23 mars et 9 avril 2020. Quel prix était le bon~? La valeur liquidative était obsolète, le
cours de l'ETF avait été poussé à la baisse, et la valeur réelle des obligations était celle que personne n'observait.

\section{Fiches de stratégie}

\begin{strategyfile}{Valeur crédit et momentum}\label{strat:s2:systematic-credit-and-bond-etf-arbitrage:factors}
\sfield{Qui paie, et pourquoi} Des investisseurs contraints par les notations et les indices de référence, qui laissent les obligations bon marché le rester et poursuivent lentement les gagnants récents.
\sfield{Instruments et plateformes} Obligations d'entreprises~; CDS comme substitut liquide.
\sfield{Signal} Spread par rapport à des pairs de notation et de duration voisines~; momentum du spread~; couverture de la duration.
\sfield{Dimensionnement et exécution} Position longue--courte ou inclinaisons par rapport à un indice de référence~; mensuelle.
\sfield{Coûts} Fourchette des obligations, large pour les émissions petites et anciennes.
\sfield{Comment elle meurt} Les coûts~; une valeur qui n'est que du risque de défaut~; les krachs du momentum aux points de retournement.
\sfield{Horizon, capacité, infrastructure} Des mois~; une base de données obligataire avec des indicateurs de cours obsolètes.
\sfield{Backtest honnête} Cours échangés, et non cours de matrice~; coûts par tranche de liquidité.
\sfield{Sources} Houweling et van Zundert (2017)~; ce chapitre~: valeur 0.46, momentum 0.14.
\end{strategyfile}

\begin{strategyfile}{Crédit à faible risque}\label{strat:s2:systematic-credit-and-bond-etf-arbitrage:lowrisk}
\sfield{Qui paie, et pourquoi} Des investisseurs qui se tournent vers le rendement en détenant de longues obligations risquées et laissent les obligations courtes et sûres bon marché par unité de risque.
\sfield{Instruments et plateformes} Obligations à courte échéance et de meilleure notation, à effet de levier ou opposées à de longues obligations risquées.
\sfield{Signal} Duration et notation.
\sfield{Dimensionnement et exécution} Risque égal de chaque côté.
\sfield{Coûts} Financement de l'effet de levier.
\sfield{Comment elle meurt} Une remontée du crédit risqué~; les coûts de financement.
\sfield{Horizon, capacité, infrastructure} Des années.
\sfield{Backtest honnête} Effet de levier aux taux de financement (funding rate) réalistes.
\sfield{Sources} Houweling et van Zundert (2017)~; ce chapitre~: 0.95.
\end{strategyfile}

\begin{strategyfile}{Arbitrage de prime et de décote d'ETF}\label{strat:s2:systematic-credit-and-bond-etf-arbitrage:etf}
\sfield{Qui paie, et pourquoi} Les détenteurs d'ETF qui vendent leurs parts contre des liquidités plus vite que les obligations ne peuvent être vendues.
\sfield{Instruments et plateformes} Parts d'ETF~; création et rachat avec l'émetteur comme participant autorisé~; les obligations.
\sfield{Signal} Le cours comparé à la valeur réelle des obligations, et non à une valeur liquidative obsolète.
\sfield{Dimensionnement et exécution} Racheter lorsque l'écart dépasse le coût de vente du panier~; paniers tels que l'émetteur les définit.
\sfield{Coûts} Vente d'obligations peu liquides en période de stress.
\sfield{Comment elle meurt} Une décote qui ne tient qu'à une valeur liquidative obsolète~; paniers modifiés pour détenir des obligations impossibles à vendre.
\sfield{Horizon, capacité, infrastructure} Des jours~; statut de participant autorisé et desks de trading obligataire.
\sfield{Backtest honnête} Estimer la valeur réelle à partir des cours échangés~; intégrer les règles de constitution des paniers.
\sfield{Sources} Haddad, Moreira et Muir (2020)~; Koont, Ma, Pastor et Zeng (2022)~; ce chapitre~: 28 jours rentables pendant le stress.
\end{strategyfile}

\begin{strategyfile}{Fourniture de liquidité sur transactions de portefeuille}\label{strat:s2:systematic-credit-and-bond-etf-arbitrage:portfolio}
\sfield{Qui paie, et pourquoi} Des clients qui paient pour négocier toute une liste d'obligations à la fois.
\sfield{Instruments et plateformes} Transactions de portefeuille cotées par rapport à l'ETF~; l'ETF comme couverture.
\sfield{Signal} Le prix de la liste comparé à l'ETF et aux cotations des obligations.
\sfield{Dimensionnement et exécution} Couvrir immédiatement avec l'ETF, puis liquider les obligations sur plusieurs jours.
\sfield{Coûts} La fourchette des obligations au fur et à mesure de leur liquidation.
\sfield{Comment elle meurt} Des listes d'obligations que seul le dealer détient~; l'ETF et les obligations qui divergent.
\sfield{Horizon, capacité, infrastructure} Des jours~; bilan.
\sfield{Backtest honnête} Exécution de chaque obligation aux cours qui étaient disponibles.
\sfield{Sources} Aucun chiffre de performance vérifié.
\end{strategyfile}

\section{Tutoriel~: quel prix est le bon}\label{tut:s2:systematic-credit-and-bond-etf-arbitrage:tut}

\begin{tutorial}
\textbf{Objectif.} Construire un univers d'obligations à structure factorielle et à cours obsolètes, négocier les \omterm{def:s2:systematic-credit-and-bond-etf-arbitrage:factor}{facteurs crédit}, et faire traverser à l'ETF et à son participant autorisé
un stress. \textbf{État final~:} le tableau et les deux figures.
\begin{tutsteps}
\item \textbf{Livres de facteurs}.
\omcode{code/firm/syscredit/firm_syscredit.py}{91}{119}{Livres valeur, momentum et faible risque, équilibrés en duration, mensuels.}\label{lst:s2:systematic-credit-and-bond-etf-arbitrage:factors}
\item \textbf{Le participant autorisé}.
\omcode{code/firm/syscredit/firm_syscredit.py}{122}{132}{Rachat à une décote supérieure au coût de vente des obligations.}\label{lst:s2:systematic-credit-and-bond-etf-arbitrage:ap}
\item \textbf{Exécuter} \texttt{factors()}, \texttt{etf()} et \texttt{fig\_syscredit.py}.
\end{tutsteps}
\textbf{Pour aller plus loin.} Estimer la valeur réelle des obligations à partir de l'ETF et des transactions récentes, négocier l'écart~; laisser le panier omettre les
obligations les moins liquides~; ajouter une annonce d'achats de la Fed qui referme la dislocation en une journée.
\end{tutorial}

\section{Construction~: crédit systématique}\label{bld:s2:systematic-credit-and-bond-etf-arbitrage:syscredit}

\begin{build}
\bfield{Rôle} Un panel d'obligations synthétiques à \omterm{def:s2:systematic-credit-and-bond-etf-arbitrage:factor}{facteurs crédit} et cours obsolètes, un ETF obligataire à valeur liquidative et cours, et arbitrage création--rachat.
\bfield{Interface} \texttt{SysCreditConfig(\dots)}, \texttt{simulate\_bonds(cfg)}, \texttt{factor\_book(sim, cfg, name)},
\texttt{ap\_arbitrage(sim, cfg)}.
\bfield{Règles} Les cours obsolètes ne se mettent à jour que les jours de négociation~; valeur liquidative calculée à partir des cours obsolètes~; arbitrage contre la valeur réelle moins le coût.
\bfield{Tests d'acceptation} \texttt{code/firm/syscredit/tests/}~: les cours obsolètes ne bougent qu'à l'occasion de transactions et la valeur liquidative en est la moyenne~; pas d'arbitrage en l'absence de
dislocation~; les livres de facteurs sont longs--courts et finis.
\bfield{Pour aller plus loin} Migration de notation~; défauts~; choix du panier.
\end{build}

\begin{omsources}
\item P.~Houweling et J.~van Zundert, ``Factor investing in the corporate bond market'', \emph{Financial Analysts Journal} 73(2), 2017.
\item V.~Haddad, A.~Moreira et T.~Muir, ``When selling becomes viral: disruptions in debt markets in the COVID-19 crisis and the Fed's response'',
NBER Working Paper 27168, 2020.
\item N.~Koont, Y.~Ma, L.~Pastor et Y.~Zeng, ``Steering a ship in illiquid waters: active management of passive funds'', NBER Working Paper 30039,
2022.
\end{omsources}

\section{Exercices}

\begin{exercise}[$\star$]\label{exo:s2:systematic-credit-and-bond-etf-arbitrage:1}
Un ETF se négocie 504 points de base sous sa valeur liquidative~; la valeur liquidative se situe 215 points de base au-dessus de la valeur réelle des obligations. De combien le prix se situe-t-il sous la valeur réelle~?
\end{exercise}

\begin{exercise}[$\star$]\label{exo:s2:systematic-credit-and-bond-etf-arbitrage:2}
Le prix se situe 300 points de base sous la valeur réelle, et vendre le panier coûte 150. Que rapporte le rachat~?
\end{exercise}

\begin{exercise}[$\star$]\label{exo:s2:systematic-credit-and-bond-etf-arbitrage:3}
Les spreads de marché s'élargissent de 150 points de base sur des obligations dont la spread duration moyenne est de 6.5. Environ combien le panier perd-il~?
\end{exercise}

\begin{exercise}[$\star\star$]\label{exo:s2:systematic-credit-and-bond-etf-arbitrage:4}
Pourquoi les \omterm{def:s2:systematic-credit-and-bond-etf-arbitrage:factor}{facteurs crédit} doivent-ils être comparés à durations voisines~?
\end{exercise}

\begin{exercise}[$\star\star$]\label{exo:s2:systematic-credit-and-bond-etf-arbitrage:5}
Pourquoi les \omterm{def:s2:systematic-credit-and-bond-etf-arbitrage:discount}{décotes d'ETF} étaient-elles plus fortes sur les obligations les plus sûres en mars 2020~?
\end{exercise}

\begin{exercise}[$\star\star$]\label{exo:s2:systematic-credit-and-bond-etf-arbitrage:6}
Pourquoi un ETF obligataire choisit-il un panier de rachat qui ne reproduit pas l'indice~?
\end{exercise}

\begin{exercise}[$\star\star\star$]\label{exo:s2:systematic-credit-and-bond-etf-arbitrage:7}
\emph{Programmation.} Relancer avec \texttt{SysCreditConfig(trade\_prob=0.05)}. Que devient la part périmée de la décote au pire moment du stress~?
\end{exercise}

\begin{exercise}[$\star\star\star$]\label{exo:s2:systematic-credit-and-bond-etf-arbitrage:8}
\emph{Trouver la faille.} «~L'ETF se traite 5~\% sous la valeur liquidative, donc l'acheter rapporte 5~\% quand la décote se résorbe.~»
\end{exercise}

\section{Problème~: Quel prix est le bon}

\begin{problem}[{Problème du week-end --- facteurs crédit et ETF obligataires}]\label{pb:s2:systematic-credit-and-bond-etf-arbitrage:1}
Les obligations synthétiques et l'ETF du chapitre, ainsi que les données publiques.

\medskip\noindent\textbf{Partie I --- Facteurs.}
\begin{enumerate}
\item Définir un \omterm{def:s2:systematic-credit-and-bond-etf-arbitrage:factor}{facteur crédit}.
\item Qu'ont trouvé Houweling et van Zundert~?
\item Décrire les obligations synthétiques et les trois signaux.
\item Donner les ratios de Sharpe et les corrélations des facteurs.
\end{enumerate}

\medskip\noindent\textbf{Partie II --- ETF.}
\begin{enumerate}[resume]
\item Définir l'arbitrage création--rachat et la \omterm{def:s2:systematic-credit-and-bond-etf-arbitrage:discount}{décote d'ETF}.
\item Qu'ont trouvé Koont, Ma, Pastor et Zeng au sujet des paniers~?
\item Comment calcule-t-on la valeur liquidative synthétique~?
\item Qu'est-ce qu'une transaction de portefeuille~?
\end{enumerate}

\medskip\noindent\textbf{Partie III --- Le stress.}
\begin{enumerate}[resume]
\item Décrire le stress introduit.
\item Décomposer la pire décote.
\item Quand le participant autorisé pourrait-il gagner de l'argent~?
\item Qu'ont trouvé Haddad, Moreira et Muir~?
\end{enumerate}

\medskip\noindent\textbf{Partie IV --- Le verdict.}
\begin{enumerate}[resume]
\item Énoncer le \emph{résultat nommé}~: les ratios de Sharpe des \omterm{def:s2:systematic-credit-and-bond-etf-arbitrage:factor}{facteurs crédit} et le rendement de l'arbitrage sur ETF dans le stress introduit.
\item Quel prix est le bon, et pourquoi personne ne peut-il le voir~?
\item Pourquoi le momentum est-il faible ici~?
\item Comment estimeriez-vous la valeur vraie dans un stress~?
\item Comment feriez-vous un backtest honnête de l'arbitrage sur ETF~?
\item Quelle fiche de stratégie exige le statut de participant autorisé~?
\item Quel lien ce chapitre entretient-il avec le chapitre~18~?
\item En une phrase~: que mesure la décote d'un ETF obligataire~?
\end{enumerate}
\end{problem}

\section{Questions d'entretien}

\begin{interviewq}[$\star$ \iqroles{trader}]\label{iq:s2:systematic-credit-and-bond-etf-arbitrage:1}
Comment la création et le rachat maintiennent-ils un ETF près de sa valeur liquidative~?
\end{interviewq}

\begin{interviewq}[$\star\star$ \iqroles{chercheur}]\label{iq:s2:systematic-credit-and-bond-etf-arbitrage:2}
Comment construiriez-vous un facteur de valeur pour les obligations d'entreprises~?
\end{interviewq}

\begin{interviewq}[$\star\star$ \iqroles{trader}]\label{iq:s2:systematic-credit-and-bond-etf-arbitrage:3}
Un client veut vendre une liste de 200 obligations d'un seul coup. Comment la pricez-vous~?
\end{interviewq}

\begin{interviewq}[$\star\star$ \iqroles{risque}]\label{iq:s2:systematic-credit-and-bond-etf-arbitrage:4}
Votre book détient des ETF obligataires contre leurs obligations. Que se passe-t-il dans un stress comme celui de mars 2020~?
\end{interviewq}

\begin{interviewq}[$\star\star$ \iqroles{développeur}]\label{iq:s2:systematic-credit-and-bond-etf-arbitrage:5}
Concevez une estimation de la juste valeur pour des obligations qui n'ont pas été échangées aujourd'hui.
\end{interviewq}

\begin{interviewq}[$\star\star\star$ \iqroles{chercheur}]\label{iq:s2:systematic-credit-and-bond-etf-arbitrage:6}
Si chaque obligation se négocie un jour sur une fraction $p$ des jours et que les valeurs vraies suivent une marche aléatoire, montrer que la valeur liquidative est en retard de $(1 - p)/p$ jours en moyenne,
et calculer l'écart périmé après un mouvement régulier de $m$ par jour.
\end{interviewq}
